\section{Introduction}
\label{sec:introduction}
The core motivation of this study is to utilize controlled neural network experiments to directly address the fundamental debate in linguistics and cognitive science regarding the nature of grammar acquisition, and to explore the robustness boundaries of computational models under degenerate input.

In classical linguistics, Chomsky (1980) proposed the "argument from poverty of the stimulus," pointing out that the real-world linguistic environment in which human children are situated is fraught with noise, such as slips of the tongue, incomplete fragments, and grammatical errors. He argued that children are able to stably acquire complex hierarchical grammar in such a degenerate environment because the brain possesses an innate "universal grammar". However, due to ethical and practical constraints, we can never conduct causal intervention experiments on human children by artificially cutting off or proportionally distorting certain types of grammatical input. Consequently, the proposition of whether innate structure is a necessary condition for grammar acquisition has long remained at the level of empirical observation and philosophical speculation, lacking rigorous empirical verification.

However, with the development of computational linguistics and deep learning, artificial neural networks provide us with an invaluable cognitive experimental platform. This is not only because of their powerful learning capabilities, but also because within neural network models, we can achieve precise control over both "innate traits" (i.e., the inductive bias of the network architecture) and the "acquired environment" (i.e., the signal-to-noise ratio of the training corpus). By injecting controlled gradients of grammatical errors into the training data of the models, we can simulate a realistic degenerate human linguistic environment, thereby rigorously testing how different internal structures react to impoverished environmental stimuli.

Furthermore, natural language is inherently hierarchically nested; yet, as Shen et al. (2019) pointed out, standard sequence models inherently lack an inductive bias for such hierarchical structure. This raises an acute mechanistic question: when correct grammatical signals in the training data are substantially diluted, how will the models' generalization strategies change? Will they stubbornly utilize the surviving correct data to maintain high-level "hierarchical rules," or will they succumb to the noise and rapidly degrade into "linear rules" that rely solely on surface word order? Quantitatively tracking this evolutionary trajectory will profoundly reveal how noise reshapes the internal syntactic representations of language models.

To deeply investigate the necessity of innate grammatical structure and to verify whether models can adhere to the "hierarchical nesting" mechanism in the presence of noise, this study constructs strict control groups based on model architectures. We compare pure sequence models lacking explicit hierarchical bias, represented by LSTMs, with models that incorporate explicit syntax tree generation mechanisms and possess strong innate structural biases, represented by RNNGs. If purely data-driven sequence models can still resist interference and maintain high-level hierarchical rules when faced with noisy corpora artificially injected with a large number of subject-verb agreement (SVA) errors, it would significantly challenge the traditional assumption that language learning in noise must rely on innate structure. Conversely, if the sequence models collapse at a specific noise threshold and rapidly degrade to a simple linear matching strategy, while the RNNG with an innate syntactic bias demonstrates strong robustness and continues to parse along a hierarchically nested path, this would inversely confirm that some form of innate structural constraint is indispensable for maintaining hierarchical computation and acquiring complex grammar in a degenerate environment. Finally, if both models collapse simultaneously at a specific noise threshold, it would further strengthen strong nativist arguments, indicating that textual input alone is insufficient to support grammar acquisition in a highly degenerate environment, and that human children must inevitably rely on richer innate mechanisms, embodied experience, or multimodal interactive input.

% OWNER: Kelly Fu (intro/conclusion per proposal)
% TODO (Kelly): The phenomenon (SVA; hierarchical rule vs. linear most-recent-noun
%       heuristic; attractors as the diagnostic case).
% TODO: The gap / our question (how does agreement behavior degrade when the
%       training signal itself is noisy?).
% TODO: Our contribution (leak-free protocol; matched LSTM vs. RNNG across five
%       noise rates; accuracy, attractor gap, surprisal).
% TODO: Roadmap.
